\subsection{Pirámide de pruebas}
Se aplica una estrategia por niveles, alineada al código real del proyecto Android. Las pruebas unitarias cubren los utilitarios de series de materias, la detección de conflictos, las heurísticas del generador de horarios, el filtrado de candidatos, la importación de horarios compartidos, la sincronización automática, el cálculo de avance académico y el componente de notificaciones de Simón.

Las pruebas de integración verifican la interacción entre casos de uso, repositorios y reglas de dominio, incluyendo sincronización, generación e importación. En la capa de interfaz se emplean pruebas instrumentadas con Jetpack Compose Testing, ActivityScenario, Hilt para pruebas y verificación de contexto Android.

Las pruebas de extremo a extremo validan los flujos críticos visibles, desde la pantalla principal hasta la generación y aplicación de horarios. Para rendimiento y carga se incorpora AndroidX Macrobenchmark con UI Automator, midiendo inicio en frío, fluidez de navegación y la traza real GenerarHorarioFlow. Las rutas completas de los módulos probados se encuentran en el Anexo \ref{ann:trazabilidad-tecnotime}.

La meta de validación fue mantener 0 fallos en las suites automatizadas y documentar los flujos manuales pendientes con instrumentos verificables. La ejecución registrada alcanzó 89 pruebas unitarias con 100\% de éxito, 7 pruebas instrumentadas con 100\% de éxito y 3 escenarios Macrobenchmark con 0 fallos.
